

\chapter{Python et pandas}\label{ch:iv3:python-and-pandas}

\begin{ivtrains}
Savoir ce que fait réellement l'interpréteur, puis ce qu'une bibliothèque fait aux données. Les énoncés précisent la version visée là où cela compte~: ici CPython 3.10, NumPy 2.2 et pandas 2.3. Les questions sur pandas portent surtout sur l'obtention du bon nombre, car un résultat rapide mais faux est l'échec habituel.
\end{ivtrains}

\begin{question}[$\star$ \iqroles{QD, Dev, QR, DS} \iqfirm{toutes} \iqlevel{junior}]\label{iq:iv3:python-and-pandas:1}
Quelles quatre lignes le \cref{lst:iv3:pp:default} affiche-t-il sous CPython, et pourquoi~? Réécrivez \texttt{add} de sorte que chaque appel sans liste parte d'une liste vide.
\omcode{code/interviews-3/02-python-and-pandas/python/pp_default.py}{2}{10}{Un argument par défaut.}\label{lst:iv3:pp:default}
\end{question}

\begin{question}[$\star$ \iqroles{QD, Dev, QR, DS} \iqfirm{toutes} \iqlevel{junior}]\label{iq:iv3:python-and-pandas:2}
Quelles deux lignes le \cref{lst:iv3:pp:alias} affiche-t-il~? Corrigez chacun des deux cas pour que chaque clé et chaque ligne possède sa propre liste. Pourquoi \texttt{dict.fromkeys(names, 0)} est-il sûr, alors que \texttt{dict.fromkeys(names, [])} ne l'est pas~?
\omcode{code/interviews-3/02-python-and-pandas/python/pp_alias.py}{1}{7}{Deux façons de partager une liste.}\label{lst:iv3:pp:alias}
\end{question}

\begin{question}[$\star$ \iqroles{QD, Dev, DS} \iqfirm{toutes} \iqlevel{junior}]\label{iq:iv3:python-and-pandas:3}
Sous CPython 3.10, quelles deux lignes le \cref{lst:iv3:pp:is} affiche-t-il~? Quelle est la différence entre \texttt{is} et \texttt{==}, et quand faut-il utiliser \texttt{is}~?
\omcode{code/interviews-3/02-python-and-pandas/python/pp_is.py}{1}{4}{Identité et égalité.}\label{lst:iv3:pp:is}
\end{question}

\begin{question}[$\star$ \iqroles{QD, Dev, DS} \iqfirm{toutes} \iqlevel{junior}]\label{iq:iv3:python-and-pandas:4}
Qu'affiche le \cref{lst:iv3:pp:keys}~? Au-delà de «~une liste peut changer et un tuple ne le peut pas~», que découle cette différence pour les clés d'un dictionnaire, et pour la mémoire~?
\omcode{code/interviews-3/02-python-and-pandas/python/pp_keys.py}{1}{7}{Ce qui peut être une clé.}\label{lst:iv3:pp:keys}
\end{question}

\begin{question}[$\star$ \iqroles{QD, Dev, DS} \iqfirm{toutes} \iqlevel{junior}]\label{iq:iv3:python-and-pandas:5}
Quelles deux lignes le \cref{lst:iv3:pp:copy} affiche-t-il~? Un collègue veut un instantané peu coûteux d'un carnet d'ordres tenu sous forme de dictionnaire de listes, pris plusieurs fois par seconde. Que lui répondez-vous~?
\omcode{code/interviews-3/02-python-and-pandas/python/pp_copy.py}{1}{9}{Deux sortes de copie.}\label{lst:iv3:pp:copy}
\end{question}

\begin{question}[$\star$ \iqroles{QD, Dev, QR, DS} \iqfirm{toutes} \iqlevel{junior}]\label{iq:iv3:python-and-pandas:6}
Quelles trois lignes le \cref{lst:iv3:pp:nan} affiche-t-il sous CPython 3.10~? Que nous apprend chacune d'elles sur le calcul du maximum d'une liste de cours pouvant contenir un NaN~?
\omcode{code/interviews-3/02-python-and-pandas/python/pp_nan.py}{1}{4}{NaN en Python pur.}\label{lst:iv3:pp:nan}
\end{question}

\begin{question}[$\star$ \iqroles{QD, Dev, DS} \iqfirm{toutes} \iqlevel{junior}]\label{iq:iv3:python-and-pandas:7}
Quelles quatre lignes le \cref{lst:iv3:pp:exhaust} affiche-t-il~? Une fonction prend un itérable de rendements, calcule sa longueur à l'aide d'une boucle, puis sa somme, et divise. Elle fonctionne sur une liste et donne un résultat faux sur un générateur, sans aucune erreur. Que renvoie-t-elle dans ce cas, et que change-t-il si la somme est calculée en premier~?
\omcode{code/interviews-3/02-python-and-pandas/python/pp_exhaust.py}{2}{8}{Un seul passage ou plusieurs.}\label{lst:iv3:pp:exhaust}
\end{question}

\begin{question}[$\star$ \iqroles{QD, Dev, QR, DS} \iqfirm{toutes} \iqlevel{junior}]\label{iq:iv3:python-and-pandas:8}
Avec NumPy 2.2, quelles trois lignes le \cref{lst:iv3:pp:overflow} affiche-t-il~? Y a-t-il un avertissement~? À quel endroit d'un pipeline de données cela poserait-il problème~?
\omcode{code/interviews-3/02-python-and-pandas/python/pp_overflow.py}{1}{7}{Entiers à largeur fixe.}\label{lst:iv3:pp:overflow}
\end{question}

\begin{question}[$\star$ \iqroles{QR, DS, QD, Analyste} \iqfirm{toutes} \iqlevel{junior}]\label{iq:iv3:python-and-pandas:9}
Dans le \cref{lst:iv3:pp:seriesand}, avec pandas 2.3, lesquelles des lignes A, B et C fonctionnent-elles, et que renvoient-elles~? Expliquez chaque échec.
\omcode{code/interviews-3/02-python-and-pandas/python/pp_series_and.py}{1}{13}{Trois façons de combiner deux conditions.\protect\hfil\protect\break}\label{lst:iv3:pp:seriesand}
\end{question}

\begin{question}[$\star\star$ \iqroles{QD, Dev, QR, DS} \iqfirm{toutes} \iqlevel{junior, confirmé}]\label{iq:iv3:python-and-pandas:10}
Sous CPython 3.10, quatre threads exécutent chacun une boucle Python pure qui fait de l'arithmétique, sur une machine disposant de quatre cœurs libres. Quelle accélération obtient-on par rapport à un seul thread, et pourquoi~? Cela rend-il \texttt{count += 1} sûr sur une variable globale partagée~? Quand les threads sont-ils utiles, et que faut-il utiliser lorsqu'ils ne le sont pas~?
\end{question}

\begin{question}[$\star\star$ \iqroles{QD, Dev, QR} \iqfirm{toutes} \iqlevel{junior}]\label{iq:iv3:python-and-pandas:11}
Écrivez un décorateur \texttt{memo} qui met en cache les résultats d'une fonction selon ses arguments. Que fait-il du nombre d'appels d'une fonction de Fibonacci récursive pour $n = 30$~? Que se passe-t-il avec une liste comme argument, et \texttt{f(1, 2)} et \texttt{f(1, b=2)} forment-ils une seule entrée ou deux~? Pourquoi ne pas écrire \texttt{cache.get(key) or f(\dots)}~?
\end{question}

\begin{question}[$\star\star$ \iqroles{QD, Dev, DS} \iqfirm{toutes} \iqlevel{junior}]\label{iq:iv3:python-and-pandas:12}
Un fichier CSV de 50 Go possède les colonnes \texttt{symbol}, \texttt{price\_ticks} et \texttt{qty}. Sur une machine disposant de 16 Go de mémoire, en Python pur, calculez le notionnel total par symbole. De quoi dépend votre consommation de mémoire~? Si un générateur garde le fichier ouvert et que son appelant cesse de le lire à mi-parcours, quand le fichier se ferme-t-il~?
\end{question}

\begin{question}[$\star\star$ \iqroles{QD, Dev} \iqfirm{toutes} \iqlevel{junior}]\label{iq:iv3:python-and-pandas:13}
Quelles deux lignes le \cref{lst:iv3:pp:classattr} affiche-t-il, et pourquoi semblent-elles incohérentes~? Comment \texttt{\_\_slots\_\_} modifierait-il cette classe, et que gagne-t-on avec un million d'instances~?
\omcode{code/interviews-3/02-python-and-pandas/python/pp_class_attr.py}{1}{15}{Un attribut de la classe.}\label{lst:iv3:pp:classattr}
\end{question}

\begin{question}[$\star\star$ \iqroles{QD, Dev} \iqfirm{toutes} \iqlevel{junior}]\label{iq:iv3:python-and-pandas:14}
Écrivez un gestionnaire de contexte, d'abord sous forme de classe puis avec \texttt{contextlib}, qui journalise «~acquire~» à l'entrée et «~release~» à la sortie, même lorsque le corps lève une exception. Que se passe-t-il si \texttt{\_\_exit\_\_} renvoie une valeur vraie~? Qu'est-ce qui ne fonctionne plus dans la version \texttt{contextlib} si l'on omet \texttt{try} et \texttt{finally}~?
\end{question}

\begin{question}[$\star\star$ \iqroles{QD, Dev} \iqfirm{toutes} \iqlevel{junior}]\label{iq:iv3:python-and-pandas:15}
Vous écrivez une dataclass \texttt{Order} dotée d'un champ \texttt{tags: list = []}. Que se passe-t-il, et quand~? Corrigez-le. Puis vous placez des objets \texttt{Order} dans un ensemble, et cela échoue aussi. Pourquoi~? Que vous apporte \texttt{frozen=True}, et que ne protège-t-il pas~?
\end{question}

\begin{question}[$\star\star$ \iqroles{QD, Dev, DS} \iqfirm{toutes} \iqlevel{junior}]\label{iq:iv3:python-and-pandas:16}
Vous disposez d'une liste de transactions, chacune étant une paire formée d'un symbole et d'un notionnel. Triez-les par symbole croissant et, à symbole égal, par notionnel décroissant. Les transactions à égalité doivent conserver leur ordre d'origine. Montrez une méthode qui fonctionne même lorsque la seconde clé ne peut pas être inversée, comme une chaîne de caractères. Dans quel ordre triez-vous~?
\end{question}

\begin{question}[$\star\star$ \iqroles{QR, DS, QD} \iqfirm{toutes} \iqlevel{junior}]\label{iq:iv3:python-and-pandas:17}
Dans le \cref{lst:iv3:pp:broadcast}, \texttt{r} contient des rendements, une ligne par jour et une colonne par action, et l'objectif est de soustraire à chaque jour la moyenne transversale. Avec NumPy 2.2, que donnent les deux affichages~? Lequel est juste, comment le voir d'un coup d'œil, et que se serait-il passé avec trois actions au lieu de deux~?
\omcode{code/interviews-3/02-python-and-pandas/python/pp_broadcast.py}{1}{5}{Soustraire la moyenne de chaque ligne.\protect\hfil\protect\break}\label{lst:iv3:pp:broadcast}
\end{question}

\begin{question}[$\star\star$ \iqroles{QR, DS, QD} \iqfirm{toutes} \iqlevel{junior}]\label{iq:iv3:python-and-pandas:18}
Qu'affiche le \cref{lst:iv3:pp:views} avec NumPy 2.2~? Lesquels parmi \texttt{b}, \texttt{c} et \texttt{d} sont des vues~? Vous conservez \texttt{huge[:10]} à partir d'un tableau de 10 Go. Que devient la mémoire~? Dans pandas 2.3, que fait \texttt{df[df.x > 0][\textquotesingle y\textquotesingle] = 1} à \texttt{df}~?
\omcode{code/interviews-3/02-python-and-pandas/python/pp_views.py}{1}{10}{Vues et copies.}\label{lst:iv3:pp:views}
\end{question}

\begin{question}[$\star\star$ \iqroles{QR, DS, QD, Analyste} \iqfirm{toutes} \iqlevel{junior}]\label{iq:iv3:python-and-pandas:19}
Un DataFrame de transactions possède les colonnes \texttt{symbol}, \texttt{day}, \texttt{price} et \texttt{qty}, avec les lignes (A, 1, 10, 100), (A, 1, 11, 300), (B, 1, 20, 50) et (A, 2, 12, 0). Calculez le prix moyen pondéré par les volumes par symbole et par jour, sans \texttt{apply}. Que donne-t-il pour chaque groupe~? Que donne \texttt{groupby(\dots).price.mean()} pour (A, 1), et pourquoi est-ce faux~?
\end{question}

\begin{question}[$\star\star$ \iqroles{QR, DS, QD, QT} \iqfirm{toutes} \iqlevel{junior}]\label{iq:iv3:python-and-pandas:20}
Pour un symbole, des transactions ont lieu aux instants entiers 1, 5 et 10, et des cotations aux instants 2, 5 et 9. Associez à chaque transaction la dernière cotation avec \texttt{pd.merge\_asof}. Quel instant de cotation obtient chaque transaction~? Que change \texttt{allow\_exact\_matches=False}, et \texttt{direction=\textquotesingle forward\textquotesingle}, et lequel de ces réglages constitue un bogue dans un backtest~? Que se passe-t-il si l'entrée n'est pas triée~?
\end{question}

\begin{question}[$\star\star$ \iqroles{QR, DS, QD, QT} \iqfirm{hedge fund, teneur de marché} \iqlevel{junior}]\label{iq:iv3:python-and-pandas:21}
Des ticks arrivent à 09:30:10 au cours de 10, à 09:30:50 à 12, à 09:31:20 à 11 et à 09:33:05 à 13. Construisez des barres d'une minute en ouverture, plus haut, plus bas et clôture avec \texttt{resample}. Que contient la barre de 09:30, et sous quelle étiquette horaire pandas la nomme-t-il par défaut~? Pourquoi cette étiquette fait-elle fuiter le futur dans un backtest, et quelle est la correction~? Que doit contenir la barre de 09:32~?
\end{question}

\begin{question}[$\star\star$ \iqroles{QR, DS, QD, Analyste} \iqfirm{toutes} \iqlevel{junior}]\label{iq:iv3:python-and-pandas:22}
Quelles trois lignes le \cref{lst:iv3:pp:align} affiche-t-il avec pandas 2.3~? Supposons que vous disposiez des cours de deux actions et qu'une action n'ait pas coté un jour~: laquelle des trois est le bogue qui ne produit aucune erreur, et que fait-il aux rendements d'une paire~?
\omcode{code/interviews-3/02-python-and-pandas/python/pp_align.py}{1}{7}{Alignement sur l'index.}\label{lst:iv3:pp:align}
\end{question}

\begin{question}[$\star\star$ \iqroles{QR, QT, DS} \iqfirm{hedge fund, teneur de marché, banque} \iqlevel{junior}]\label{iq:iv3:python-and-pandas:23}
Un backtest tire des rendements quotidiens normaux indépendants, de moyenne nulle, pose \texttt{signal = np.sign(ret)} et calcule \texttt{pnl = signal * ret}. Quel ratio de Sharpe annualisé (moyenne divisée par l'écart type, multipliée par $\sqrt{252}$) obtient-on, pour un long échantillon~? Démontrez-le. Quelle est la correction en un mot dans pandas, et que devient le Sharpe~? Comment rendre ce type de bogue difficile à écrire dans une bibliothèque de recherche~?
\end{question}

\begin{question}[$\star\star$ \iqroles{QR, DS, QD, Analyste} \iqfirm{toutes} \iqlevel{junior}]\label{iq:iv3:python-and-pandas:24}
Quelles quatre lignes affiche le \cref{lst:iv3:pp:nanpandas} avec pandas 2.3~? Un collègue exécute \texttt{prices.fillna(0)} avant de calculer les rendements. Qu'est-ce qui ne va pas, et que feriez-vous à la place~?
\omcode{code/interviews-3/02-python-and-pandas/python/pp_nan_pandas.py}{1}{8}{NaN dans pandas.}\label{lst:iv3:pp:nanpandas}
\end{question}

\begin{question}[$\star\star$ \iqroles{QD, DS, QR} \iqfirm{toutes} \iqlevel{junior, confirmé}]\label{iq:iv3:python-and-pandas:25}
Une table compte 100 millions de lignes~: \texttt{symbol}, une chaîne de caractères à 4\,000 valeurs distinctes, \texttt{price}, un flottant, et \texttt{qty}, un entier dont l'étendue dépend des données. Avec pandas 2.3, quels types choisiriez-vous pour économiser de la mémoire~? Environ combien prend chaque colonne~? Un collègue propose \texttt{float32} pour les prix et pour les cumuls de volumes. Montrez ce que cela casse.
\end{question}

\begin{question}[$\star\star$ \iqroles{QD, Dev, QR} \iqfirm{toutes} \iqlevel{junior}]\label{iq:iv3:python-and-pandas:26}
Écrivez des tests pytest pour une fonction qui calcule un prix moyen pondéré par les volumes par groupe, en utilisant \texttt{parametrize}. Pourquoi \texttt{0.1 + 0.2 == 0.3} échoue-t-il, et que permet \texttt{pytest.approx} par défaut~? \texttt{1e-13 == pytest.approx(0)} passe-t-il, et \texttt{1e-9 == pytest.approx(0)}~? Quels cas limites ajoutez-vous~?
\end{question}

\begin{question}[$\star\star$ \iqroles{QD, Dev, QR} \iqfirm{toutes} \iqlevel{junior}]\label{iq:iv3:python-and-pandas:27}
Quelles quatre lignes affiche le \cref{lst:iv3:pp:round} sous CPython 3.10~? Une commission doit être arrondie au centime, les moitiés étant arrondies vers le haut, et la valeur vaut exactement 2.675. Comment procédez-vous~?
\omcode{code/interviews-3/02-python-and-pandas/python/pp_round.py}{1}{6}{Arrondi.}\label{lst:iv3:pp:round}
\end{question}

\begin{question}[$\star\star$ \iqroles{QR, QD, LLE} \iqfirm{toutes} \iqlevel{junior, confirmé}]\label{iq:iv3:python-and-pandas:28}
Sous CPython 3.10, que renvoient \texttt{sum([1e16, 1.0, -1e16])} et \texttt{math.fsum([1e16, 1.0, -1e16])}, et pourquoi~? Écrivez votre propre sommation compensée, en un seul passage et à mémoire constante. La méthode classique de Kahan donne-t-elle le bon résultat pour cette liste~?
\end{question}

\begin{question}[$\star\star$ \iqroles{QD, Dev, LLE} \iqfirm{crypto, teneur de marché} \iqlevel{junior, confirmé}]\label{iq:iv3:python-and-pandas:29}
Un processus \texttt{asyncio} possède une coroutine heartbeat qui doit émettre toutes les 10~ms, et un gestionnaire qui effectue parfois 200~ms de travail bloquant à l'intérieur d'une coroutine. Que devient le heartbeat, et pourquoi~? Corrigez cela pour une entrée ou sortie bloquante, puis pour un calcul pur. Que se passe-t-il si l'on appelle une coroutine et qu'on oublie le \texttt{await}~?
\end{question}

\begin{question}[$\star\star\star$ \iqroles{QR, DS, QD} \iqfirm{toutes} \iqlevel{junior, confirmé}]\label{iq:iv3:python-and-pandas:30}
Une table contient 10 millions de lignes de \texttt{symbol}, \texttt{x} (un rendement du facteur) et \texttt{y} (le rendement de l'action), pour 5\,000 symboles, avec des NaN dans l'une ou l'autre colonne. Dans pandas, sans \texttt{apply} et sans boucle sur les symboles, calculez la pente des moindres carrés de \texttt{y} sur \texttt{x} par symbole. Pour un symbole dont les lignes sont (1, 2), (2, 4), (3, 6) et (100, NaN), quelle est la bonne pente~? Que donne-t-on si l'on centre chaque colonne séparément puis que l'on somme, en laissant pandas ignorer les NaN~? Que doit renvoyer un symbole ayant une seule ligne~?
\end{question}

\begin{question}[$\star\star\star$ \iqroles{QR, DS, QD, QT} \iqfirm{hedge fund, teneur de marché} \iqlevel{junior, confirmé}]\label{iq:iv3:python-and-pandas:31}
Une table de cotations contient \texttt{t}, \texttt{symbol}, \texttt{bid} et \texttt{ask}, sur plusieurs jours. Certaines lignes ont un côté nul ou un marché croisé, certaines partagent un horodatage, certaines arrivent dans le désordre, et un mid de 1\,000 se trouve parmi des mids de 100. Dans pandas, écrivez une étape de nettoyage, puis les rendements logarithmiques à une minute du mid, par symbole, sans rendement à cheval sur une nuit. Dans quel ordre s'enchaînent les étapes~? Pourquoi un test de valeurs aberrantes ne doit-il pas utiliser une fenêtre centrée~? Si l'on rejette un mid s'écartant de plus de 5~\% de la médiane des 20 mids qui le précèdent, que se passe-t-il pour un vrai saut de 6~\%~?
\end{question}

\begin{question}[$\star\star\star$ \iqroles{QD, QR, LLE, QT} \iqfirm{toutes} \iqlevel{junior, confirmé}]\label{iq:iv3:python-and-pandas:32}
Les valeurs arrivent une à une. Écrivez une moyenne exponentiellement pondérée à mémoire constante par flux, qui reproduit exactement \texttt{ewm(alpha=a).mean()} de pandas 2.3, NaN compris. Pour $a = 0.1$ et les valeurs 1, 2, 3, que donne-t-elle à la troisième valeur, et que donne \texttt{adjust=False}~? Quand les deux coïncident-elles~? Maintenant, les ticks arrivent à des instants irréguliers~: modifiez le calcul pour qu'un poids soit divisé par deux à chaque période de temps fixe. Que fait votre version lorsque deux ticks partagent un horodatage, et après un écart d'un jour~?
\end{question}

\begin{question}[$\star\star\star$ \iqroles{QR, QD, Risque} \iqfirm{toutes} \iqlevel{junior, confirmé}]\label{iq:iv3:python-and-pandas:33}
Avec NumPy, tirez $n$ échantillons de $k$ normales centrées réduites de matrice de corrélation $C$. Un collègue calcule \texttt{L = np.linalg.cholesky(C)}, tire \texttt{z} avec un échantillon par ligne, et renvoie \texttt{z @ L}. Pour deux variables de corrélation 0.8, quelle matrice de covariance les échantillons possèdent-ils~? Donnez la bonne ligne et démontrez-la. Que se passe-t-il avec des corrélations deux à deux de 0.9, 0.9 et $-0.9$, pourquoi aucune donnée ne peut-elle les produire, et que faites-vous~? Que se passe-t-il avec une corrélation égale à exactement 1~?
\end{question}

\begin{question}[$\star\star\star$ \iqroles{QR, QD, Risque} \iqfirm{banque, hedge fund, teneur de marché} \iqlevel{junior, confirmé}]\label{iq:iv3:python-and-pandas:34}
Dans le modèle de Black--Scholes, prenez un spot et un strike égaux à 100, un taux de 5~\%, une volatilité de 20~\% et une maturité d'un an~: le call vaut 10.4506. Pricez-le par Monte-Carlo en NumPy, avec des paires antithétiques et une erreur type, de sorte que le nombre de trajectoires ne soit pas limité par la mémoire. Que se passe-t-il si l'on traite les deux membres de chaque paire comme indépendants lors du calcul de l'erreur~? L'erreur est-elle alors trop grande ou trop petite, et cela dépend-il du payoff~? Donnez un contrôle exact de votre code qui ne nécessite pas de prix d'option en forme fermée.
\end{question}

\begin{question}[$\star\star\star$ \iqroles{QD, QT, QR} \iqfirm{toutes} \iqlevel{junior, confirmé}]\label{iq:iv3:python-and-pandas:35}
Des ordres signés arrivent, et chacun est tronqué pour que la position reste dans $\pm L$~; l'ordre suivant part de la position tronquée. Un collègue calcule la position par \texttt{np.clip(np.cumsum(orders), -L, L)}. Donnez le plus petit exemple possible où cela est faux. Écrivez la version correcte. La version à deux bornes peut-elle être vectorisée~? Si seule une borne supérieure s'applique, existe-t-il une forme close avec \texttt{cumsum}~?
\end{question}

\begin{question}[$\star\star\star$ \iqroles{QD, Dev, LLE} \iqfirm{toutes} \iqlevel{junior, confirmé}]\label{iq:iv3:python-and-pandas:36}
Sous CPython 3.10, un service construit un nouveau \texttt{Pricer} à chaque requête, comme dans le \cref{lst:iv3:pp:growth}. La mémoire croît et ne redescend jamais, et \texttt{gc.collect()} ne libère rien. Qu'est-ce qui retient les objets~? Après 50 requêtes, que renvoie \texttt{Pricer.price.cache\_info()}, et pourquoi le cache est-il inutile autant que nuisible~? Comment repérer ce type de croissance dans un processus en production, et qu'est-ce que votre outil ne voit pas~? Corrigez cela. \texttt{maxsize=128} est-il une correction~?
\omcode{code/interviews-3/02-python-and-pandas/python/iv_growth.py}{2}{15}{Un cache sur une méthode.}\label{lst:iv3:pp:growth}
\end{question}

\begin{question}[$\star\star\star$ \iqroles{QR, QD, DS} \iqfirm{toutes} \iqlevel{junior, confirmé}]\label{iq:iv3:python-and-pandas:37}
Un backtest fixe sa graine aléatoire et lit les mêmes données, pourtant deux exécutions affichent des nombres différents. Une forme réduite figure dans le \cref{lst:iv3:pp:hashseed}. Sous CPython 3.10, pourquoi deux exécutions de ce programme peuvent-elles afficher des lignes différentes, dans les deux instructions d'affichage~? Le même phénomène se produirait-il avec des codes de symboles entiers~? Comment corrigez-vous cela, et pourquoi fixer \texttt{PYTHONHASHSEED} est-il la mauvaise correction~? Quelles autres sources de ce problème vérifieriez-vous~?
\omcode{code/interviews-3/02-python-and-pandas/python/pp_hashseed.py}{1}{5}{Deux exécutions, deux réponses.}\label{lst:iv3:pp:hashseed}
\end{question}

\begin{question}[$\star\star\star\star$ \iqroles{QR, DS, QD} \iqfirm{hedge fund} \iqlevel{confirmé}]\label{iq:iv3:python-and-pandas:38}
Une table contient \texttt{date}, \texttt{symbol}, \texttt{price} et \texttt{signal}, pour 3\,000 actions sur 2\,500 jours, avec des NaN dans \texttt{signal}. Pour chaque date, calculez la corrélation de rang de Spearman entre le signal et le rendement du lendemain de chaque action, sans \texttt{apply} et sans boucle sur les dates. Pourquoi le rendement suivant ne doit-il pas provenir de \texttt{df.price.shift(-1)} sur la table triée~? Si le signal et le rendement suivant sont conjointement normaux de corrélation 0.1, quelle moyenne attendez-vous~? À quel point est bruitée une valeur quotidienne avec 500 actions~? Que change-t-on dans la statistique de Student si la cible devient un rendement sur cinq jours~?
\end{question}

\begin{question}[$\star\star\star\star$ \iqroles{QD, LLE, DS} \iqfirm{toutes} \iqlevel{confirmé}]\label{iq:iv3:python-and-pandas:39}
En NumPy seulement, écrivez une jointure as-of. Vous disposez des tableaux \texttt{trade\_sym}, \texttt{trade\_t}, \texttt{quote\_sym} et \texttt{quote\_t}~: les symboles sont des codes entiers de 0 à 4\,999, les temps sont des \texttt{int64} en nanosecondes depuis 1970, proches de $1.8 \times 10^{18}$, et rien n'est trié. Pour chaque transaction, renvoyez la ligne de la cotation la plus récente du même symbole à l'instant ou avant son temps, ou $-1$. Pour les cotations de même temps et de même symbole, la dernière arrivée l'emporte. Aucune boucle sur les symboles n'est permise. Pourquoi la clé \texttt{sym * 10**19 + t}, ou tout multiplicateur assez grand pour les temps, échoue-t-elle~? Que coûte votre version~?
\end{question}

\begin{question}[$\star\star\star\star$ \iqroles{QR, QD, LLE} \iqfirm{toutes} \iqlevel{confirmé}]\label{iq:iv3:python-and-pandas:40}
\texttt{x} et \texttt{y} sont des tableaux NumPy de rendements quotidiens d'environ 1~\%, avec des NaN. Pour chaque ligne, calculez la pente de \texttt{y} sur \texttt{x} sur les $w$ dernières lignes, en n'utilisant que les lignes où les deux valeurs sont présentes, ou NaN si moins de \texttt{min\_obs} lignes remplissent la condition. N'utilisez aucune boucle Python sur les lignes, et un temps en $O(n)$. La méthode habituelle différencie une somme cumulée. Un mauvais tick met alors $10^9$ dans \texttt{x} à la ligne 10 sur 1\,000. Avec $w = 50$, que donne cette méthode à la ligne 999, où la vraie pente vaut près de 2~? Peut-on éviter les dégâts tout en restant en $O(n)$~? Que montre \texttt{rolling(50).var()} de pandas 2.3.3 à cette ligne~?
\end{question}
